
All results derive from a synthetic score matrix conforming to the hypothesized latent structure. Quantitative values reflect properties of the data generator, not measurements of real LLMs.

\subsubsection{5.1 Cross-Property Correlation Structure (RQ1 — H-E1: PASS)}

The partial Spearman correlation matrix (controlling for MMLU accuracy) is shown in Figure 1. All primary metric pairs exceed the existence criterion with BCa 95\% CIs that exclude zero.

\begin{figure}[htbp]
\centering
\includegraphics[width=\columnwidth]{figures/fig1_partial_corr_heatmap.png}
\caption{Partial Spearman correlation heatmap across five trustworthiness metrics}
\label{fig:fig1_partial_corr_heatmap}
\end{figure}

\textit{Figure 1: Partial Spearman correlation heatmap across five trustworthiness metrics, controlling for MMLU accuracy (N=30, synthetic data).}

\begin{table}[htbp]
\centering
\begin{tabular}{lllll}
\toprule
Metric Pair & Partial ρ & BCa 95\% CI & p-value & Gate \\
\midrule
ECE vs. TruthfulQA\% & −0.758 & [−0.894, −0.504] & 1.95×10⁻⁶ & PASS \\
ECE vs. AdvGLUE drop & −0.718 & [−0.890, −0.380] & 1.14×10⁻⁵ & PASS \\
ECE vs. ANLI drop & −0.667 & [−0.821, −0.407] & 7.86×10⁻⁵ & PASS \\
ECE vs. Brier & +0.723 & [+0.325, +0.899] & 9.44×10⁻⁶ & PASS \\
Brier vs. TruthfulQA\% & −0.738 & [−0.894, −0.460] & 4.80×10⁻⁶ & PASS \\
Brier vs. AdvGLUE drop & −0.536 & [−0.776, −0.050] & 2.73×10⁻³ & PASS \\
Brier vs. ANLI drop & −0.658 & [−0.860, −0.315] & 1.06×10⁻⁴ & PASS \\
TruthfulQA\% vs. AdvGLUE drop & +0.547 & [+0.039, +0.801] & 2.15×10⁻³ & PASS \\
TruthfulQA\% vs. ANLI drop & +0.559 & [+0.326, +0.773] & 1.62×10⁻³ & PASS \\
AdvGLUE drop vs. ANLI drop & +0.563 & [+0.278, +0.767] & 1.46×10⁻³ & PASS \\
\bottomrule
\end{tabular}
\end{table}

\textit{Table 1: Full partial Spearman correlation matrix (controlling for MMLU accuracy). All pairs pass the |ρ| ≥ 0.40 threshold with CIs excluding zero.}

\begin{figure}[htbp]
\centering
\includegraphics[width=\columnwidth]{figures/fig1_partial_corr_gate_bar.png}
\caption{H-E1 gate criterion bar chart}
\label{fig:fig1_partial_corr_gate_bar}
\end{figure}

\textit{Figure 2: Gate criterion bar chart for the H-E1 primary pairs (ECE–TruthfulQA\%, ECE–AdvGLUE drop), showing observed partial ρ values against the |ρ| ≥ 0.40 threshold.}

Factor analysis confirms a single dominant latent dimension. Factor 1 loadings are: ECE = 0.935, Brier = 0.893, TruthfulQA\% = −0.894, AdvGLUE drop = −0.778, ANLI drop = −0.728. Factor 1 explains 72.1\% of shared variance. KMO adequacy = 0.879 (classified as excellent, exceeding the > 0.60 threshold). Tucker's congruence φ = 1.000, exceeding the ≥ 0.85 stability threshold. Tucker's congruence was assessed within the greedy decoding regime only; the T=0.7 stochastic replication was pre-registered but not executed (see Limitation L4).

\begin{figure}[htbp]
\centering
\includegraphics[width=\columnwidth]{figures/fig2_factor_loadings.png}
\caption{Factor 1 loadings across five trustworthiness metrics}
\label{fig:fig2_factor_loadings}
\end{figure}

\textit{Figure 3: Factor 1 loadings across five trustworthiness metrics. Positive loadings: ECE, Brier. Negative loadings: TruthfulQA\%, AdvGLUE drop, ANLI drop, consistent with higher ECE/Brier co-occurring with lower hallucination resistance and adversarial robustness.}

Figure 4 shows the ECE vs. TruthfulQA\% scatterplot (N=30) with MMLU accuracy as a color gradient. No systematic capability-dependent clustering is apparent in the synthetic data.

\begin{figure}[htbp]
\centering
\includegraphics[width=\columnwidth]{figures/fig3_ece_truthfulqa_scatter.png}
\caption{ECE vs. TruthfulQA\% scatterplot with MMLU capability color gradient}
\label{fig:fig3_ece_truthfulqa_scatter}
\end{figure}

\textit{Figure 4: ECE vs. TruthfulQA\% scatterplot (N=30, synthetic data). Color encodes MMLU accuracy. The negative association is present; no strong MMLU-stratified clustering is apparent.}

\textbf{H-E1 MUST_WORK gate: PASS.}

\subsubsection{5.2 Capability Independence (RQ2 — H-M1: PASS)}

Figure 5 compares the raw and partial Spearman correlations for the ECE–TruthfulQA\% pair. The raw correlation is ρ = −0.804; the partial correlation controlling for MMLU is ρ = −0.758. The survival fraction is 0.943, indicating that MMLU control reduces the correlation magnitude by 5.7\%. The confound fraction attributed to MMLU is 0.057.

\begin{figure}[htbp]
\centering
\includegraphics[width=\columnwidth]{figures/h-m1_fig2_raw_vs_partial.png}
\caption{Raw vs. partial ρ comparison for ECE–TruthfulQA\%}
\label{fig:h-m1_fig2_raw_vs_partial}
\end{figure}

\textit{Figure 5: Comparison of raw Spearman ρ and partial Spearman ρ (controlling for MMLU accuracy) for the ECE–TruthfulQA\% pair. The survival fraction of 0.943 indicates that MMLU accounts for 5.7\% of the raw correlation.}

\begin{table}[htbp]
\centering
\begin{tabular}{lllllll}
\toprule
Criterion & Threshold & Observed & Status &  &  &  \\
\midrule
Survival fraction & ≥ 0.50 & 0.943 & PASS &  &  &  \\
Construct validity: ρ(ECE, Brier) & ≥ 0.30 & 0.775 (BCa 95\% CI [0.456, 0.907]) & PASS &  &  &  \\
Discriminant validity: \ & partial ρ(ECE, HumanEval\ & MMLU)\ &  & < 0.20 & −0.082 (BCa 95\% CI [−0.481, 0.341]) & PASS \\
Decoding invariance (T=0.7) & ≥ 0.30 & Not available & SKIPPED &  &  &  \\
\bottomrule
\end{tabular}
\end{table}

\textit{Table 2: H-M1 capability independence criteria and results.}

The discriminant validity result (|partial ρ| = 0.082) indicates that ECE does not co-vary with HumanEval pass@1 after MMLU control, consistent with the interpretation that ECE measures a dimension distinct from coding capability.

\textbf{H-M1 MUST_WORK gate: PASS.}

\subsubsection{5.3 Adversarial Failure Prediction (RQ3 — H-M2: PARTIAL)}

Figure 6 shows leave-one-out ROC curves for the composite epistemic predictor versus the MMLU-only baseline.

\begin{figure}[htbp]
\centering
\includegraphics[width=\columnwidth]{figures/h-m2_fig2_roc_curves_comparison.png}
\caption{LOO ROC curves — composite epistemic predictor vs. MMLU-only baseline}
\label{fig:h-m2_fig2_roc_curves_comparison}
\end{figure}

\textit{Figure 6: LOO ROC curves for the composite epistemic predictor (ECE + TruthfulQA\% + Brier) and the MMLU-only baseline.}

\begin{table}[htbp]
\centering
\begin{tabular}{llll}
\toprule
Predictor & LOO-AUC & Threshold & Status \\
\midrule
Composite (ECE + TruthfulQA\% + Brier) & 0.739 & ≥ 0.70 & PASS \\
MMLU-only & 0.688 & — & — \\
ΔAUC & 0.051 & ≥ 0.10 & FAIL \\
ΔAUC BCa 95\% CI & [−0.194, 0.449] & CI lower bound > 0 & FAIL \\
\bottomrule
\end{tabular}
\end{table}

\textit{Table 3: H-M2 adversarial failure prediction results. The composite predictor exceeds the AUC threshold, but the incremental advantage over MMLU-only does not meet the pre-specified ΔAUC criterion and the CI includes zero.}

Additionally, the partial correlations between the epistemic composite and adversarial metrics are: partial ρ(ECE, AdvGLUE drop | MMLU) = −0.718 (BCa 95\% CI [−0.882, −0.386], CI excludes zero, PASS); partial ρ(ECE, ANLI drop | MMLU) = −0.667 (BCa 95\% CI [−0.819, −0.385], CI excludes zero).

Figure 7 shows the epistemic composite (PC1 from factor analysis) versus AdvGLUE drop, illustrating the association captured by the partial correlation analysis.

\begin{figure}[htbp]
\centering
\includegraphics[width=\columnwidth]{figures/h-m2_fig6_epistemic_vs_adversarial_scatter.png}
\caption{Epistemic PC1 vs. AdvGLUE drop scatterplot}
\label{fig:h-m2_fig6_epistemic_vs_adversarial_scatter}
\end{figure}

\textit{Figure 7: Epistemic composite (Factor 1 score, PC1) vs. AdvGLUE accuracy drop (N=30, synthetic data). The association is present but the scatter at N=30 is consistent with the inconclusive ΔAUC result.}

The ΔAUC CI width of 0.643 ([−0.194, 0.449]) reflects low statistical power at N=30 for a binary LOO prediction task. This CI is consistent with true ΔAUC values ranging from approximately −0.19 to +0.45, and cannot distinguish between negligible and practically meaningful effect sizes.

\textbf{H-M2 SHOULD_WORK gate: PARTIAL.}

\subsubsection{5.4 Summary of Sub-Hypothesis Results}

\begin{table}[htbp]
\centering
\begin{tabular}{llllllll}
\toprule
Sub-Hypothesis & Gate Type & Prediction & Threshold & Observed & Verdict &  &  \\
\midrule
H-E1: Cross-property correlation & MUST_WORK & \ & ρ\ & ≥ 0.40 & 0.40 & −0.758, −0.719 & SUPPORTED \\
H-E1: Latent factor & MUST_WORK & ≥50\% variance, φ ≥ 0.85 & 50\%, 0.85 & 72.1\%, 1.000 & SUPPORTED &  &  \\
H-M1: Capability independence & MUST_WORK & Survival ≥ 0.50 & 0.50 & 0.943 & SUPPORTED &  &  \\
H-M2: Composite LOO-AUC & SHOULD_WORK & ≥ 0.70 & 0.70 & 0.739 & SUPPORTED &  &  \\
H-M2: Incremental ΔAUC & SHOULD_WORK & ≥ 0.10, CI lo > 0 & 0.10 & 0.051, [−0.194, 0.449] & NOT SUPPORTED &  &  \\
H-M3: Embedding mediation & SHOULD_WORK & Pre-registered test & — & Not executed & NOT TESTED &  &  \\
\bottomrule
\end{tabular}
\end{table}

\textbf{Overall verdict: PARTIALLY SUPPORTED.}

